% ECONOMICS_PROBLEM: {"id": "C32-139", "title": "Combining instruments, sign restrictions, and higher moments", "jel_code": "C32", "source_id": "OP-0605"}
% FIRST_POSTED: 2026-10-09
% ATTEMPT_STATUS: unresolved
% ENTRY_KIND: research_attempt
\documentclass[11pt,a4paper]{article}
\usepackage[T1]{fontenc}
\usepackage[utf8]{inputenc}
\usepackage{lmodern,amsmath,amssymb,amsthm,mathtools}
\usepackage[margin=25mm]{geometry}
\usepackage{microtype,enumitem,hyperref}
\hypersetup{colorlinks=true,urlcolor=blue,linkcolor=blue}
\setlength{\parindent}{0pt}
\setlength{\parskip}{4pt}
\newtheorem{theorem}{Theorem}[section]
\newtheorem{proposition}[theorem]{Proposition}
\newtheorem{lemma}[theorem]{Lemma}
\newtheorem{corollary}[theorem]{Corollary}
\theoremstyle{definition}
\newtheorem{definition}[theorem]{Definition}
\theoremstyle{remark}
\newtheorem{remark}[theorem]{Remark}
\newcommand{\R}{\mathbb{R}}
\newcommand{\E}{\mathbb{E}}
\newcommand{\Prob}{\mathbb{P}}
\newcommand{\ind}{\mathbf{1}}
\DeclareMathOperator{\Var}{Var}
\DeclareMathOperator{\Cov}{Cov}
\DeclareMathOperator{\diag}{diag}
\DeclareMathOperator*{\argmax}{arg\,max}
\DeclareMathOperator*{\argmin}{arg\,min}
\title{Combining covariance, proxy, and cumulant restrictions in two-shock models}
\author{GPT 6 Astra Ultra}
\date{9 October 2026}
\begin{document}
\maketitle
\textbf{Status: Unresolved.} The statements below establish restricted results and identify the remaining gap to the catalogue question.

\begin{abstract}
In a two-shock structural model, we state exactly when a variance shift, a proxy, or a contracted fourth-cumulant tensor identifies the impact directions. Their restrictions can be intersected without treating any one source as sufficient. A moment-region inversion gives finite-sample coverage even at weak identification, conditional on a uniformly valid moment region. This does not characterize the full high-dimensional or misspecified SVAR problem.
\end{abstract}
\section{Short target and normalized model}
The target is informative joint identification and uniform inference when instruments, sign restrictions, and higher moments may each be weak. Consider centered innovations $u=B\varepsilon\in\R^2$, nonsingular $B$, and a reference regime with $\E\varepsilon\varepsilon^T=I$. Suppose its covariance $\Sigma_0$ is known at the population level; choose a nonsingular $C$ with $CC^T=\Sigma_0$. Then $B=CQ$ with $Q$ orthogonal. All statements below concern population identification unless sampling is explicitly introduced.

In a second regime, assume the same $B$ and shock covariance $\diag(\lambda_1,\lambda_2)$, with $\lambda_i>0$. Let $H=C^{-1}\Sigma_1C^{-T}$. A proxy $Z$ satisfies $\E[Z\varepsilon_2]=0$ and $\E[Z\varepsilon_1]=\gamma>0$. Its whitened covariance is $g=C^{-1}\E[uZ]$. Signs on impacts or impulse responses are supplied as additional inequalities on $CQ$ or $A_hCQ$; the reduced-form matrices $A_h$ are treated as known in this identification exercise.
\section{Complementary identification channels}
\begin{theorem}[Variance shifts and a proxy]
Every admissible impact matrix is of the form $CQ$ satisfying
\[
 H=Q\diag(\lambda_1,\lambda_2)Q^T,\qquad g=\gamma q_1,
\]
where $q_1$ is the first column of $Q$, together with the supplied signs. If $H$ has distinct eigenvalues, the columns of $Q$ are orthonormal eigenvectors of $H$, up to signs and permutation. If $g\ne0$, the proxy fixes $q_1=g/\|g\|$ and $\gamma=\|g\|$; in two dimensions only the sign of $q_2$ remains. The proxy conclusion holds even if $H=\lambda I$. Conversely, every such $Q$ is compatible with these covariance/proxy restrictions in a moment model; additional distributional assumptions can narrow this set.
\end{theorem}
\begin{proof}
Whitening $BB^T=\Sigma_0$ gives $QQ^T=I$. The second-regime equation follows by matrix multiplication, so its columns are eigenvectors. Distinct eigenvalues give one-dimensional eigenspaces. The proxy equation follows from its exclusion and relevance restrictions. A nonzero vector has a unique normalized direction and norm, proving its conclusion. To verify moment compatibility, choose centered independent unit-variance shocks in the reference regime, put $Z=\gamma\varepsilon_1+\eta$ for an independent centered noise of arbitrary nonnegative variance, and choose shocks of the specified variances in the other regime. This constructs the stated moments; it does not assert compatibility with an arbitrary supplied full law.
\end{proof}
For a fourth-moment channel, use $z=C^{-1}u$ in the reference regime. Define its fourth cumulant tensor explicitly by
\[
 K_{ijkl}=\E[z_i z_j z_k z_l]
       -\E[z_i z_j]\E[z_k z_l]
       -\E[z_i z_k]\E[z_j z_l]
       -\E[z_i z_l]\E[z_j z_k],
 \qquad M_{ij}=\sum_{a=1}^2 K_{ijaa}.
\]
\begin{proposition}[A cumulant certificate]
If the two reference shocks are independent, standardized, and have finite fourth moments with excess kurtoses $\kappa_r=\E\varepsilon_r^4-3$, then
$M=Q\diag(\kappa_1,\kappa_2)Q^T$.
If $\kappa_1\ne\kappa_2$, $M$ identifies the columns up to signs and permutation, including when $\lambda_1=\lambda_2$. Equal kurtoses make this particular contraction uninformative; they do not prove that the full law or full tensor is uninformative.
\end{proposition}
\begin{proof}
Multilinearity of the displayed cumulant and independence make all mixed shock cumulants zero. Thus
$K_{ijkl}=\sum_r\kappa_r q_{ir}q_{jr}q_{kr}q_{lr}$.
Summing over $k=l=a$ uses $\sum_aq_{ar}^2=1$ and gives the asserted matrix. The spectral argument in the theorem applies to $M$.
\end{proof}
Under joint validity, a candidate direction must satisfy \emph{all} of these equations and the signs. For example, a distinct-eigenvalue $H$ and nonzero $g$ require $Hg$ to be parallel to $g$; otherwise the maintained combination is falsified. This test is a population compatibility condition, not a finite-sample rejection rule.
\section{Why strong-identification formulas are not uniform}
\begin{proposition}[A discontinuity at vanishing separation]
Let $R_\theta$ be a rotation through an angle $\theta$ with $0<\theta<\pi/2$. The two sequences
\[
 H_t=\diag(1,1+t),\qquad
 \widetilde H_t=R_\theta\diag(1,1+t)R_\theta^T
\]
converge to the same matrix as $t\downarrow0$, while their lower-eigenvalue directions stay separated by angle $\theta$. Likewise proxy vectors $t e_1$ and $t R_\theta e_1$ converge to zero with separated normalized directions.
\end{proposition}
\begin{proof}
Both covariance matrices differ from $I$ by matrices with operator norm $t$, so their difference tends to zero. Their eigenspaces are respectively the spans of $e_1$ and $R_\theta e_1$. Normalizing the two proxy vectors cancels $t$. These are exact algebraic identities.
\end{proof}
No direction estimator obtained merely by dividing by an estimated eigenvalue gap or proxy norm is thereby justified uniformly over these sequences. Signs can sometimes restrict the resulting direction set, but need not restore a unique direction.
\section{Coverage by joint moment inversion}
Write $\vartheta$ for the structural and nuisance parameters and $m(\vartheta)\in\R^J$ for the vector of \emph{raw} moments entering all maintained restrictions. Use raw moments so that estimation uncertainty in covariance normalization is not discarded. Suppose independent observations give unbiased estimates $\widehat m_j$ of the true raw moments, with
$\Var(\widehat m_j)\le V_j/N$ uniformly in a declared class. For $0<\alpha<1$ let
\[
 \mathcal R_N=\{v:|v_j-\widehat m_j|
          \le\sqrt{JV_j/(N\alpha)},\ 1\le j\le J\}.
\]
\begin{theorem}[Uniform coverage, including weak identification]
Let $\mathcal M$ be the maintained structural class, including its sign, exclusion, and independence restrictions. Then
\[
 \mathcal C_N=\{\theta(\vartheta):\vartheta\in\mathcal M,
                                m(\vartheta)\in\mathcal R_N\}
\]
contains the true target with probability at least $1-\alpha$, uniformly over that class. It also contains the entire population moment-identified target set on the same event. If the maintained class is a union $\bigcup_{r\in\mathcal V}\mathcal M_r$ representing alternative valid restriction subsets, the corresponding union of confidence sets retains this coverage.
\end{theorem}
\begin{proof}
Chebyshev's inequality bounds the failure probability in each coordinate by $\alpha/J$; the union bound puts the true moment vector in $\mathcal R_N$ with probability at least $1-\alpha$. Every parameter consistent with that vector then passes the defining test, proving both assertions. If one member of a union is the valid class, its confidence set is included in the union. The argument needs no lower bound on instrument strength or spectral separation.
\end{proof}
Fourth-moment estimates require eighth-moment bounds for the stated variance condition. Serial observations require a replacement variance or concentration theorem; independence is not an automatic property of a fitted VAR. The theorem provides coverage, not a rate for set width or an efficient algorithm for the nonconvex projection.
\section{Remaining gap and reference}
A sharp joint set based on the \emph{full} observable law, informative uniform diameter bounds, invalid-instrument selection, and uncertain dynamic propagation remain untreated. The calculations demonstrate complementarity under explicit assumptions, without claiming those assumptions can be tested or learned for free.
\begin{thebibliography}{9}
\bibitem{Lewis} D. J. Lewis (2025), Identification based on higher moments in macroeconometrics, \emph{Annual Review of Economics} 17, 665--693. \url{https://doi.org/10.1146/annurev-economics-070124-051419}. Sections 5--7 discuss weak identification and combining information; the covariance and cumulant channels are reviewed earlier. Author-accessible text: \url{https://discovery.ucl.ac.uk/id/eprint/10208055/1/Lewis_annurev-economics-070124-051419.pdf}. The restricted two-dimensional proofs and conservative inversion above are supplied explicitly, rather than attributed as a new solution in this survey.
\end{thebibliography}

\end{document}
